
\def\PUDcodigo{ACE015}

\if\COMPILAR0
   \def\PUDcodacad{04507.38}
\else
   \def\PUDcodacad{04506.37}
\fi

\def\PUDdisciplina{Empreendedorismo}

\def\PUDsemestre{S8}

\def\PUDhoras{40}

%NOVOS ---------------------------------------------------
\def\PUDhorasTeoria{40}
\def\PUDhorasPratica{0}

\if\COMPILAR0
   \def\PUDinterECA{\hyperref[PUD:ACE017]{Introdução à Engenharia (04507.5)}}
\else
   \def\PUDinterECA{\hyperref[PUD:ACE017]{Introdução à Engenharia (04506.5)}}
\fi

\def\PUDatuaisECA{---}

\def\PUDrecursos{Projetor multimídia; Quadro branco e pincel.}

%----------------------------------------------------------

\def\PUDcreditos{2}

\def\PUDprerequisito{---}

\def\PUDpreacad{---}

\def\PUDobjetivos{Capacitar-se para o desenvolvimento das habilidades empreendedoras através de atividades teóricas e práticas; Fazer uso das tecnologias da informação, adequando-as aos novos modelos organizacionais e dos processos e sistemas de inovação tecnológica; Selecionar ideias e pesquisar necessidades de mercado; Definir critérios para avaliação do potencial de um novo negócio e dos recursos necessários para desenvolvê-lo e implementá-lo. Articular competências gerais do curso para construção na implementação de um plano de negócios. Conhecer os fundamentos da administração com enfoque a engenharia.}

\def\PUDementa{Aspectos relacionados à prática do empreendedorismo.  Gerenciando recursos empresariais.  Plano de 
negócios: importância, estrutura e apresentação.  Caminhos a seguir e recursos disponíveis para o 
empreendedor.  A gestão empreendedora e suas implicações para as organizações.  O papel e a importância do comportamento empreendedor nas organizações.  O perfil dos profissionais empreendedores nas organizações.  Processos grupais e coletivos, processos de autoconhecimento, autodesenvolvimento, criatividade, comunicação e liderança.  A iniciativa e tomada de decisão. FUndamentos de Administração para engenharia. }

\def\PUDconteudo{ &  Empreendedorismo: O mundo globalizado e seus desafios e potencialidades;  Conhecendo o empreendedorismo (introdução, estudos, definições de diversos autores);  Características dos empreendedores;  Competências e Habilidades: persistência, comprometimento, exigência de qualidade e eficiência, persuasão e rede de contatos, independência e autoconfiança, busca de oportunidades, busca de informações, planejamento e monitoramento sistemático, estabelecimento de metas, correr riscos 
calculados;  Identificação de oportunidades de negócio.  
\\ 
 &  Gerenciando os recursos empresariais: Gerenciando a equipe;  Gerenciando a produção;  Gerenciando o marketing;  Gerenciando as finanças 
\\ 
 &  Plano de negócios: A importância do plano de negócios;  Estrutura do plano de negócios;  Elementos de um plano de negócios eficiente;  Exemplo de um plano de negócios.  
\\ 
 &  Assessoria para o negócio: Buscando assessoria: incubadoras de empresas, SEBRAE, Franchising, Universidades e institutos de pesquisa, assessoria jurídica e contábil;  Criando a empresa;  Questões legais de constituição da empresa: tributos, marcas e patentes;  Apresentação de planos de negócios.
 \\
 & Fundamentos de Administração. Enfoque sistêmico. A construção de uma teoria administrativa, com foco no aumento da produtividade; papéis do gerente. Administração da qualidade nos EUA e modelo japonês. Abordagem humanística e teoria comportamental. Estilos de administração. Visão sociotécnica, grupos semiautônomos de trabalho. Administração participativa, administração por resultados. Administração no presente. }

\def\PUDmetodo{Aulas expositórias que podem ser teóricas e/ou práticas, onde as práticas no laboratório serão marcadas no decorrer da disciplina.}

\def\PUDavalia{A avaliação será feita com aplicação de provas teóricas e/ou práticas, além da possibilidade de inclusão de trabalhos e seminários no decorrer da disciplina.}

\def\PUDbibbasica{ &   \href{www.biblioteca.ifce.edu.br/index.asp?codigo_sophia=62049}{Drucker}, P. F.  Inovação e espírito empreendedor (entrepreneurship): prática e princípios.  São Paulo, SP: Cengage Learning, 2015.  378p.  ISBN: 9788522108596.
%\\ 
 %&   MAXIMIANO, A.  C.  A.  Administração para empreendedores: fundamentos da criação e da gestão de novos negócios.  São Paulo: Prentice-Hall, 2006.  ISBN 9788576050889
\\
 & \href{www.biblioteca.ifce.edu.br/index.asp?codigo_sophia=61750}{Maximiano}, Antonio Cesar Amaru.  Administração para empreendedores.  2º ed.  São Paulo, SP: Pearson Prentice Hall, 2011.  240p.  ISBN: 9788576058762.
%\\ 
 %&   DORNELAS, José C.  A.  Empreendedorismo na prática: mitos e verdades do empreendedores de sucesso.  Rio de Janeiro: Elsevier, 2007.  ISBN 9788535227611
\\ 
 &   \href{www.biblioteca.ifce.edu.br/index.asp?codigo_sophia=65131}{Dornelas}, José C. A.  Empreendedorismo: transformando idéias em negócios.  6ª ed.  São Paulo, SP: Atlas, 2016.  267p.  ISBN: 9788597003932. }

\def\PUDbibcomplementar{ & \href{www.biblioteca.ifce.edu.br/index.asp?codigo_sophia=61985}{Chiavenato}, Idalberto.  Empreendedorismo: dando asas ao espírito empreendedor.  4ª ed.  Barueri, SP: Manole, 2012.  315p.  ISBN: 9788520432778.
\\
&   \href{www.biblioteca.ifce.edu.br/index.asp?codigo_sophia=62402}{Ferreira}, Ademir A.  Gestão empresarial: de Taylor aos nossos dias.  São Paulo, SP: Cengage Learning, 2016.  247p.  ISBN: 9788522100985.
\\ 
 &   \href{www.biblioteca.ifce.edu.br/index.asp?codigo_sophia=65379}{Salim}, C. S.  Introdução ao empreendedorismo: despertando a atitude empreendedora.  Rio de Janeiro, RJ: Elsevier, 2010.  245p.  ISBN: 9788535234664.
\\
 &   \href{www.biblioteca.ifce.edu.br/index.asp?codigo_sophia=55920}{Martinelli}, Dante Pinheiro; Joyal, André.  Desenvolvimento Local e o Papel das Pequenas e Médias Empresas.  E-book.  Manole.  356p.  ISBN: 9788520416662.
\\
 &   \href{www.biblioteca.ifce.edu.br/index.asp?codigo_sophia=24197}{Gauthier}, Fernando Álvaro Ostuni.  Empreendedorismo.  Curitiba, PR: Livro Técnico, 2010.  120p.  ISBN: 9788563687173. }

\def\PUDautor{Fabrício Bandeira da Silva}

\def\PUDdata{2012-04-25}

\def\PUDrevisao{1}

\def\PUDrevisor{Samuel Vieira Dias}

\def\PUDdatarevisao{2017-05-23}

\PUDcorpo
